% TOP_PROBLEM: {"id": 6900009, "title": "Configuration Spaces of Tensegrities — Problem 9", "category_id": 6, "category": {"id": 6, "name": "geometry", "display_name": "Geometry", "description": "Euclidean and non-Euclidean geometry, geometric structures.", "slug": "geometry", "order_index": 6, "created_at": "2026-07-31T15:26:25.671Z"}, "status": "open"}
% FIRST_POSTED: null
% ENTRY_KIND: statement_only
% Imported from: batch18.tex; UnsolvedMath ID 6900009
\documentclass[11pt]{article}
\usepackage[margin=1in]{geometry}
\usepackage{amsmath,amssymb,mathtools}
\usepackage{fontspec,unicode-math}
\setmainfont{Latin Modern Roman}
\setsansfont{Latin Modern Sans}
\setmathfont{Latin Modern Math}
\usepackage{xurl,hyperref}
\hypersetup{colorlinks=true,urlcolor=blue,linkcolor=blue}
\newcommand{\sref}[2]{\hyperlink{source:#1:#2}{[S#2]}}
\newcommand{\cataloguescope}[1]{\paragraph{Recorded mathematical question.}#1}
\begin{document}
% UnsolvedMath ID 6900009; source catalogue code AMR-068-0009
\setcounter{section}{6900008}
\section{Configuration Spaces of Tensegrities — Problem 9}\label{6900009}
{\small\sffamily UnsolvedMath ID 6900009 · Geometry}
\subsection{Definitions and mathematical statement}

\paragraph{Shared notation.}$\mathbb N=\{1,2,\ldots\}$, and $\mathbb Z,\mathbb Q,\mathbb R,\mathbb C$ denote the integers, rationals, reals and complex numbers. Any other conventions are specified locally.


Let $G$ be a simple graph with labeled vertices $1,\ldots,n$. A configuration is $P=(p_1,\ldots,p_n)\in(\mathbb R^d)^n$; coincident vertices are allowed. A self-stress is a real edge assignment $w_{ij}=w_{ji}$, set to zero on missing edges, satisfying equilibrium at every vertex:
\[\sum_{j\ne i}w_{ij}(p_j-p_i)=0\quad(1\leq i\leq n).\]
Let $W(G,P)$ be the vector space of all such self-stresses. Negative and positive edge values are the source’s cable and strut signs. The base $B_d(G)$ is the full labeled configuration space $(\mathbb R^d)^n$, stratified by the following fiber equivalence: $W(G,P)$ and $W(G,Q)$ are equivalent if a homeomorphism between them preserves the sign of every edge coordinate of every stress. A stratum is a maximal connected set of configurations with equivalent fibers; these strata are semialgebraic. This base is not quotiented by Euclidean or affine motions. Write $K_n$ for the complete graph on $n$ labeled vertices. In the planar model, the ordered geometric variables lie in the real projective plane $\mathbb{RP}^2=\mathbb R^2\cup\ell_\infty$, which adds a line of directions so that parallel affine lines meet at a point at infinity. Elementary conditions are coincidence of two points, collinearity of three points, and the five-point intersection relation $p_i=[p_j,p_{j\prime};p_k,p_{k\prime}]$. The last is the closure in $(\mathbb{RP}^2)^5$ of the relation $p_i=(p_jp_{j\prime})\cap(p_kp_{k\prime})$ where the two joining lines are well-defined and distinct. This closed relation also covers degenerate endpoint pairs where a joining line is undefined. For example, when both joining lines coincide, $p_i$ may be any point on their common line. A finite system may introduce auxiliary points in $\mathbb{RP}^2$; its conditional solutions among the original points are those for which suitable projective auxiliary points exist. For a graph $G$ and $k\geq1$, the source calls $\{P:\dim W(G,P)\geq k\}$ a $(G,k)$-stratum (possibly a union of base strata). It is geometric if it is a finite union of conditional-solution sets of such geometric systems. These planar descriptions motivate the higher-dimensional question.
\paragraph{Statement recorded in the supplied source.}
Develop a theory of geometric conditions describing tensegrity strata in dimensions $d\geq3$; in particular, find elementary geometric conditions in three and higher dimensions analogous to the planar incidence conditions and their conditional systems. \sref{6900009}{2}
\subsection{Short English statement}
Develop higher-dimensional geometric incidence descriptions of tensegrity strata, extending the planar theory.
\subsection{Sources}
\begin{description}
\item[\hypertarget{source:6900009:1}{S1}] ULAM.AI and the UnsolvedMath Contributors, \emph{UnsolvedMath: A Curated Collection of Open Mathematics Problems}, record 6900009 (AMR-068-0009). CC BY 4.0. \url{https://huggingface.co/datasets/ulamai/UnsolvedMath}.
\item[\hypertarget{source:6900009:2}{S2}] Franck Doray, Oleg Karpenkov and Jan Schepers, \emph{Geometry of configuration spaces of tensegrities}, arXiv:0806.4976v2 (2008), Sections 2 and 5.2--5.3, Definitions 2.1 and 2.6, Theorem 2.8, Definitions 5.1 and 5.4 and Problem 3, pp. 3--6 and 17--18. \url{https://arxiv.org/pdf/0806.4976v2}.
\end{description}
\label{end:6900009}
\end{document}
